\section{From Everlasting Death, Deliver Us}

Hell is the work of a refusal, not of a decree that predestines. ``God predestines no one to go to hell,'' the \emph{Catechism} teaches; ``for this, a willful turning away from God (a mortal sin) is necessary, and persistence in it until the end.''\endnote{\emph{Catechism of the Catholic Church} 1037.} The Second Council of Orange said it with unusual vehemence: that some have been predestined to evil by divine power, \latin{non solum non credimus, sed etiam, si sunt qui tantum mali credere velint, cum omni detestatione illis anathema dicimus}---we not only do not believe it, but if there are any who wish to believe so great an evil, we say anathema to them with all detestation.\endnote{Second Council of Orange (529), conclusion, DS 397; cited by the \emph{Catechism} at 1037 together with Trent, session VI, canon 17, DS 1567.} The God who threatens hell is the God who swore by his own life that he does not desire it for anyone: ``As I live, saith the Lord God, I desire not the death of the wicked, but that the wicked turn from his way, and live. Turn ye, turn ye from your evil ways: and why will you die, O house of Israel?''\endnote{Ezechiel 33:11, Douay--Rheims (Challoner).} ``The Lord\ldots{} dealeth patiently for your sake, not willing that any should perish, but that all should return to penance.''\endnote{2 Peter 3:9, Douay--Rheims (Challoner).}

\subsection{The door that stands open}

The saints spoke about hell to the living, for whom the door of mercy is still open. Cyprian, who had just set the eternal fire before the persecutor Demetrianus, turned in the next breath to offer him the remedy:

\begin{quote}
Provide, therefore, while you may, for your safety and your life\ldots{} Nor let any one be restrained either by his sins or by his years from coming to obtain salvation. To him who still remains in this world no repentance is too late. The approach to God's mercy is open, and the access is easy to those who seek and apprehend the truth.\endnote{Cyprian, \emph{To Demetrianus} 25.}
\end{quote}

\noindent Christ himself promises: ``All that the Father giveth to me shall come to me: and him that cometh to me, I will not cast out.''\endnote{John 6:37, Douay--Rheims (Challoner).} The thief on the cross came in the last hour and was received. Judas, as the Church's prayer says, received the punishment of his crime; not because his sin was too great to be forgiven, but because he would not ask. No sin committed in this life is beyond the reach of the blood of Christ and the absolution of his Church. The only sin that cannot be forgiven is the one that refuses forgiveness to the end.\endnote{\emph{Catechism of the Catholic Church} 982, 1864; cf.\ Matthew 12:31--32.}

The saints name the way of escape plainly. It is confession: the humble telling of sins that Tertullian set against the fire, and that Alphonsus's lost soul wished he had made every week. It is perseverance: the daily return to God, the avoidance of occasions, the prayer for the grace to persevere that Alphonsus places on the lips of his reader again and again, ``obtain for me holy perseverance.''\endnote{Alphonsus, \emph{Preparation for Death}, Consideration XXVIII, second point, Affections and Prayers, p.~345.} And it is prayer, the Church's prayer above all, in which she has always asked for exactly this deliverance.

\subsection{The Church's prayer}

In the heart of the Roman Canon the priest extends his hands over the offering and prays \latin{diesque nostros in tua pace disponas, atque ab æterna damnatione nos eripi, et in Electorum tuorum jubeas grege numerari}, in the words the \emph{Catechism} gives: ``Grant us your peace in this life, save us from final damnation, and count us among those you have chosen.''\endnote{Roman Canon, \latin{Hanc igitur}; Latin of the Canon as printed in \emph{The Roman Missal, Translated into the English Language for the Use of the Laity} (Philadelphia: Eugene Cummiskey, 1861); English from \emph{Catechism of the Catholic Church} 1037, quoting the Roman Missal, Eucharistic Prayer I.} Bede records the tradition that it was Gregory the Great who added these words to the Canon: ``he added three phrases full of great goodness and perfection: `And dispose our days in thy peace, and preserve us from eternal damnation, and rank us in the number of thy elect, through Christ our Lord.'\,''\endnote{Bede, \emph{Ecclesiastical History of the English Nation} II.1 (London: J. M. Dent; New York: E. P. Dutton, 1910), Internet Medieval Sourcebook, accessed 16 October 2026.} The pope whom Bede credits with this petition is the same who preached that many are called and few chosen, and who asked his people who could know how few of them were numbered in the flock of God's elect.

On All Souls' Day, when the Church offers the Holy Sacrifice for all her dead, she prays at the offertory, ``O Lord Jesus Christ, King of glory, deliver the souls of all the faithful departed from the pains of hell and from the deep pit; deliver them from the lion's mouth, that hell engulf them not, nor they fall into darkness, but let Michael, the holy standard-bearer, bring them into the holy light which Thou once didst promise to Abraham and his seed.''\endnote{Mass of All Souls' Day, Offertory, in Lasance and Walsh, \emph{The New Roman Missal} (1945).} At the absolution of the dead she sings, ``Deliver me, O Lord, from eternal death on that dreadful day when the heavens and the earth shall be moved, and Thou shalt come to judge the world by fire.''\endnote{The Absolution after the Mass for the Dead, responsory \latin{Libera me, Domine, de morte aeterna}, in Lasance and Walsh, \emph{The New Roman Missal} (1945).} In the sequence of the Mass for the Dead the soul pleads before the Judge in its own name:

\begin{quote}
\latin{Confutatis maledictis,\\
Flammis acribus addictis,\\
Voca me cum benedictis.}

\medskip
While the wicked are confounded,\\
Doom'd to flames of woe unbounded,\\
Call me with Thy Saints surrounded.\endnote{\latin{Dies irae}, stanza on \latin{Confutatis maledictis}; Latin in Lasance and Walsh, \emph{The New Roman Missal} (1945), Masses for the Dead; English by William J. Irons in \emph{Hymns Ancient and Modern}, complete edition (London, 1890), no.~398, stanza 8. Stanzas of the sequence close the second volume of this series, \emph{What Death Does Not Finish}.}
\end{quote}

\noindent And in the Litany of the Saints, which the Church sings on Holy Saturday, the invocations of the saints give way to the cry of the whole people: ``From all evil,'' and the answer, ``O Lord, deliver us''; ``From all sin''; ``From everlasting death.''\endnote{Litany of the Saints, in Lasance and Walsh, \emph{The New Roman Missal} (1945), Holy Saturday; the response ``O Lord, deliver us'' is printed with the first of these petitions and serves for those that follow.}

The Church prays for this deliverance because the danger is real, and she prays with confidence because she prays to the Saviour of whom it is written, ``For God sent not his Son into the world, to judge the world: but that the world may be saved by him.''\endnote{John 3:17, Douay--Rheims (Challoner).}

\subsection{Resolution}

Francis de Sales ends his meditation on hell with a few simple acts, and the reader may end this consideration with them. First, confess before God that you have deserved it. Then, ``Resolve henceforth to act differently, and to rescue yourself from this abyss. Resolve on distinct definite acts by which you may avoid sin, and thereby eternal death.'' And last: ``Give thanks, offer yourself, pray.''\endnote{Francis de Sales, \emph{Introduction to the Devout Life} I.15, Affections and Resolutions, p.~42.}

The year is ending, and the Church will soon begin it again. On the first Sunday of Advent the 1962 Missal reads once more of the signs in the sun and moon and stars and of the Son of man coming with great power and majesty, and the Church hears him say, ``when these things begin to come to pass, look up and lift up your heads, because your redemption is at hand.''\endnote{Luke 21:25--28, Douay--Rheims (Challoner); Luke 21:25--33 is the Gospel of the First Sunday of Advent in the 1962 Missal, and Luke 21:25--28, 34--36 in Year C of the Postconciliar Missal.} The Judge who will say \emph{Depart} to those on his left will say \emph{Come} to those on his right, and the Church sets hell before her children so that they may hear the second word.
